% appendices/codebook.tex
% Expert Survey Codebook -- companion to appendices/survey-questionnaire.tex

\chapter{Expert Survey Codebook}
\label{ch:codebook}

This codebook documents the data collected through the expert survey on
generative-AI--driven disinformation. It complements the survey instrument
reproduced in \texttt{appendices/survey-questionnaire.tex} and is intended to
enable independent reproduction of the empirical analyses reported in the
associated publications.

\section{Dataset Overview}
\label{sec:codebook-overview}

\begin{table}[htbp]
\centering
\caption{Expert survey dataset overview.}
\label{tab:codebook-overview}
\begin{tabular}{lll}
\toprule
\textbf{Property} & \textbf{Value} & \textbf{Notes} \\
\midrule
File         & \texttt{expert\_survey\_responses.csv} & Google Forms export \\
Respondents  & 54                                     & After consent filter \\
Items        & 54                                     & 5 sections (see below) \\
Period       & July 2025 -- March 2026                & Rolling recruitment \\
Languages    & English                                & Single-language instrument \\
\bottomrule
\end{tabular}
\end{table}

\noindent\textbf{Anonymization.}
The released dataset retains only pseudonymized respondent identifiers.
Respondents who explicitly consented to attribution at the end of the survey
are listed by name in the ``Acknowledgements'' field; their open-ended
responses may be quoted with attribution. All other identifying fields
(name, affiliation, profile link) are stripped from the public release.

\noindent\textbf{Recruitment.}
Targeted invitation of disinformation researchers, AI ethicists, journalists,
fact-checkers, and policy experts via professional mailing lists, snowball
sampling, and direct outreach. No financial incentive was offered.

\section{Variable Groups}
\label{sec:codebook-variables}

The 54 items group into five thematic blocks, summarized in
Table~\ref{tab:codebook-blocks}.

\begin{table}[htbp]
\centering
\caption{Survey blocks and item counts.}
\label{tab:codebook-blocks}
\begin{tabular}{llrl}
\toprule
\textbf{Block} & \textbf{Section} & \textbf{Items} & \textbf{Type} \\
\midrule
Screening      & 1                & 2  & Categorical / link  \\
Demographics   & 2                & 4  & Ordinal / categorical \\
Threats        & 3                & 20 & 7-point Likert (4 mod.\ $\times$ 5 contexts) \\
Risks          & 4                & 2  & Multi-select + free text \\
Mitigation     & 5                & 11 & Likert + forced ranking \\
Open feedback  & 6--9             & 15 & Free text \\
\bottomrule
\end{tabular}
\end{table}

\subsection{Threat Perceptions (20 items)}
Respondents rated four AI-generated content modalities
(\textit{Text}, \textit{Images}, \textit{Audio / voice cloning},
\textit{Video / deepfakes}) on a 7-point Likert scale
($1 = $ Not a threat, $7 = $ Severe threat) for the general case
and within four specific domains
(\textit{Political}, \textit{Health}, \textit{Financial}, \textit{Social}).
Each domain block also includes one optional free-text follow-up.

\subsection{Mitigation Effectiveness (11 items)}
Five mitigation strategies were rated on the same 7-point scale,
then forced-ranked from 1 (most effective) to 5 (least effective):
\begin{enumerate}[label=(\alph*), leftmargin=2.5em]
    \item Digital watermarking / provenance standards
    \item Platform enforcement policies
    \item Technical detection tools
    \item Government regulation (e.g., AI~Act, DSA)
    \item Media / public literacy programmes
\end{enumerate}
Two single-choice items capture the perceived MOST and LEAST effective
strategy independent of the ranking, plus one free-text item on the
biggest adoption obstacle.

\subsection{Urgent Risks (2 items)}
A multi-select item with 12 pre-specified risk categories asked respondents
to select up to three urgent risks. A follow-up free-text item invited
justification or real-world incident references.

\section{Descriptive Summary}
\label{sec:codebook-descriptives}

Tables~\ref{tab:desc-threats} and~\ref{tab:desc-mitigation} report
mean / SD across all 54 respondents on the headline Likert items.
Full per-group breakdowns are provided in the main paper.

\begin{table}[htbp]
\centering
\caption{Threat perceptions (1--7 Likert, $N=54$).}
\label{tab:desc-threats}
\begin{tabular}{lcccc}
\toprule
              & \textbf{Text} & \textbf{Images} & \textbf{Audio} & \textbf{Video} \\
\midrule
General       & 5.57 (1.36) & 5.74 (1.13) & 5.91 (1.25) & 6.19 (1.19) \\
Political     & 5.76 (1.41) & 5.98 (1.18) & 5.96 (1.20) & 6.31 (1.15) \\
Health        & 5.80 (1.32) & 5.13 (1.54) & 5.02 (1.66) & 5.13 (1.55) \\
Financial     & 5.30 (1.51) & 4.91 (1.44) & 5.65 (1.62) & 5.31 (1.46) \\
Social        & 5.59 (1.35) & 5.76 (1.20) & 5.65 (1.44) & 6.00 (1.29) \\
\bottomrule
\end{tabular}
\end{table}

\begin{table}[htbp]
\centering
\caption{Mitigation strategies: effectiveness ratings (1--7) and forced rank
(1 = most, 5 = least; $N=54$).}
\label{tab:desc-mitigation}
\begin{tabular}{lcc}
\toprule
\textbf{Strategy} & \textbf{Rating} & \textbf{Rank} \\
\midrule
Digital watermarking      & 5.20 (1.52) & 2.98 (1.43) \\
Platform enforcement      & 4.91 (1.40) & 3.26 (1.40) \\
Technical detection       & 4.57 (1.38) & 3.28 (1.27) \\
Government regulation     & 5.24 (1.63) & 3.07 (1.51) \\
Media / public literacy   & 5.07 (1.65) & 2.98 (1.48) \\
\bottomrule
\end{tabular}
\end{table}

\noindent\textbf{Top urgent risks (votes among 54 respondents, top 5):}
Election interference via deepfake video~(42),
Cyberattacks / phishing~(25),
Harassment / non-consensual deepfake porn~(25),
Audio-clone financial fraud~(21),
Large-scale text spambots~(18).

\section{Data Access and Reuse}
\label{sec:codebook-access}

The de-identified response file, aggregated tables, and the survey instrument
are archived at Harvard Dataverse
(\texttt{doi:10.7910/DVN/BXO2QA}, open, CC0~1.0). The fuller response data
underlying the Wave~1 analysis is deposited at Zenodo
(\texttt{doi:10.5281/zenodo.18703601}); because that deposit retains material
that cannot be released openly, access is granted on request for academic
research. The survey instrument is reproduced in full in
\texttt{appendices/survey-questionnaire.tex}.

This repository is licensed under the MIT License; see \texttt{LICENSE}.
Please cite the publications listed in \texttt{CITATION.cff} when reusing
the instrument or the data.
